% =============================================================================
% DA-GPS GNN, QSTS non-uniqueness, and physics-informed training (GNN2)
% Compile:
%   pdflatex da_gps_physics_nonuniqueness_report.tex
%   bibtex da_gps_physics_nonuniqueness_report
%   pdflatex da_gps_physics_nonuniqueness_report.tex
%   pdflatex da_gps_physics_nonuniqueness_report.tex
% =============================================================================
\documentclass[11pt,a4paper]{article}

\usepackage[margin=1in]{geometry}
\usepackage{amsmath,amssymb}
\usepackage{bm}
\usepackage{booktabs}
\usepackage{hyperref}
\usepackage{graphicx}

\hypersetup{colorlinks=true, linkcolor=black, citecolor=black, urlcolor=black}

\title{DA-GPS Surrogate Modeling for Quasi-Static Time-Series\\
with Autonomous Voltage Controls:\\
Non-Uniqueness, Closed-Loop Physics, and Evaluation on the IEEE~8500 Feeder}
\author{GNN2 technical report (draft)}
\date{\today}

\begin{document}
\maketitle

\begin{abstract}
Quasi-static time-series (QSTS) analysis of distribution feeders with autonomous
voltage controllers does not admit a one-to-one map from exogenous injections to
converged voltages and device states: deadbands, delays, and hysteresis make the
same load and PV inputs consistent with multiple valid coupled solutions unless
controller history is tracked~\cite{deboever2017qsts}.
This report summarizes the problem framing, the DA-GPS graph-neural surrogate
(multitask voltage, regulator, capacitor, and meta-auxiliary heads), optional
physics-informed nodal loss, and evaluation against OpenDSS native daily QSTS on
the large unbalanced IEEE~8500-node test feeder with 2~MW preinstalled PV.
Representative notebook runs report monitor-node voltage MAE $\approx 4.3\times
10^{-3}$~p.u.\ and deployment wall-time speedups of roughly $3$--$4\times$
relative to sequential OpenDSS \texttt{Solve()} calls on the same display grid.
\end{abstract}

% =============================================================================
\section{Introduction and Problem Statement}
\label{sec:problem}
% =============================================================================

Distribution interconnection and hosting-capacity studies increasingly require
QSTS simulation to capture fast PV variability, on-load tap changer (LTC)
actions, switched capacitor banks, and smart-inverter Volt--VAr
curves~\cite{deboever2017qsts,reno2017vq,ieee1547}.
On the IEEE~8500-node unbalanced OpenDSS model (\texttt{Master-PV2MW-inv.dss}),
three device classes operate autonomously: \texttt{PV2} Volt--VAr control, nine
switched single-phase capacitor banks, and twelve single-phase LTC
regulators~\cite{gnn2_ieee8500_avc,ieee8500}.

\subsection{Non-uniqueness of injection-to-output mappings}

At each 5-minute snapshot, exogenous inputs are load and irradiance multipliers
applied before \texttt{Solve()}:
\begin{equation}
  \mathbf{u}(t) = \big[m_L(t),\; m_{\mathrm{irr}}(t),\; t\big]^\top .
  \label{eq:inputs}
\end{equation}
Let $\boldsymbol{\kappa}(t)$ collect LTC taps, capacitor switch states, and
post-control inverter reactive power.  A converged OpenDSS state satisfies
\begin{equation}
  \mathbf{v}(t) = f_{\mathrm{DSS}}\!\big(\mathbf{u}(t),\boldsymbol{\kappa}(t)\big),
\end{equation}
but $\boldsymbol{\kappa}$ is \emph{not} a deterministic function of
$\mathbf{u}(t)$ alone.

Sandia and Georgia Tech document the root cause in distribution QSTS with
autonomous controls~\cite{deboever2017qsts}:
\begin{quote}
\emph{Because the deadbands in controllers create a hysteresis in the state of
the system, approximating control logic models can become extremely complex
without modeling the actual logic of the controllers.  Models based on power
injections cannot be used when controllers are considered since multiple discrete
system states would be valid for the same power injections.  For example, a
machine learning model, which seeks to establish a one to one mapping between
the QSTS inputs and outputs, is not able to learn the correlation because the
same inputs will yield multiple valid possible power flow solutions.}
\end{quote}
The report lists \emph{multiple valid power flow solutions} among the principal
barriers to fast QSTS~\cite{deboever2017qsts}.  Subsequent Sandia work on
rapid QSTS via machine learning therefore samples representative time windows,
predicts aggregate indices (e.g.\ yearly tap changes), or tracks discrete system
states rather than assuming a static injection-only
map~\cite{blakely2018evaluation,reno2017pvsc,reno2017vq,blakely2018isgt}.

\paragraph{Empirical demonstration (GNN2).}
The \texttt{nonunique\_four\_scenario\_demo} experiment in \texttt{nonunique.ipynb}
(cell~0) holds $\mathbf{u}(t)$ fixed and varies only initial controller seeds
(\texttt{init\_low}, \texttt{init\_high}, \texttt{init\_default}) or uses cold
snapshot recompiles without carry-forward (\texttt{snapshot\_cold}).  Divergent
tap, capacitor, and voltage trajectories under identical injections illustrate
path dependence and non-uniqueness on the IEEE~8500 model.

% =============================================================================
\section{Closed-Loop Power Flow with Autonomous Controls}
\label{sec:closedloop}
% =============================================================================

OpenDSS does not treat controllers as post-processing corrections.  At each time
step it alternates ac power-flow solves with control updates until device laws are
satisfied or \texttt{maxcontroliter} is reached~\cite{opendss,gnn2_ieee8500_avc}:
\begin{equation}
  (\mathbf{x}^\star,\boldsymbol{\kappa}^\star)=
  \find_{\mathbf{x},\boldsymbol{\kappa}}
  \Big\{
    \mathrm{PF}(\mathbf{x},\boldsymbol{\kappa};\mathbf{u}(t))=\mathbf{0},\;
    \mathcal{C}(\mathbf{x},\boldsymbol{\kappa})=\mathbf{0}
  \Big\},
  \label{eq:coupled}
\end{equation}
where $\mathrm{PF}(\cdot)=\mathbf{0}$ is nodal ac balance and
$\mathcal{C}(\cdot)=\mathbf{0}$ collects Volt--VAr curves, capacitor hysteresis
with voltage override and switching delay, and LTC deadband relay
logic~\cite{gnn2_ieee8500_avc}.

\paragraph{Interpretation.}
Solving distribution power flow with autonomous controllers is a \emph{closed-loop
control system}: the plant is the unbalanced network, the controller state
$\boldsymbol{\kappa}$ has memory (hysteresis, delays, warm-started taps), and the
fixed point \eqref{eq:coupled} couples algebraic PF constraints with hybrid
discrete--continuous device laws.  Snapshot-only surrogates that ignore
$\boldsymbol{\kappa}(t{-}1)$ therefore target the wrong map for daily QSTS.

Native OpenDSS daily mode (\texttt{mode=daily}) marches sequentially with
warm-start carry-forward of voltages and controller states; the GNN2 daily-compare
scripts compile once and call \texttt{Solve()} for each display step on that
path~\cite{opendss}.

% =============================================================================
\section{Approach: DA-GPS Multitask GNN}
\label{sec:approach}
% =============================================================================

\subsection{Architecture and targets}

The DA-GPS (Distribution-Aggregate Graph Perceiver with System tokens) surrogate
is implemented in \texttt{train\_da\_gps\_multitask\_complex\_voltage\_gine.py}.
It uses GINE message passing on the MV-aggregated IEEE~8500 graph, learnable
system tokens with cross-attention, and multitask heads:
\begin{itemize}
  \item complex voltage magnitude/angle (Re/Im head, denormalized for metrics);
  \item per-device capacitor bank states (BCE);
  \item per-regulator tap positions (cross-entropy on discrete tap classes);
  \item meta-auxiliary scalars (e.g.\ \texttt{PV2} post-solve $P/Q$, feeder
        losses).
\end{itemize}
Training data come from chunked OpenDSS QSTS snapshots on the unbalanced
8500-node feeder (thousands of scenarios); positional encodings and static edge
attributes supply topology that flat MLP baselines omit.

\paragraph{Why multitask helps under non-uniqueness.}
Rather than learning $\mathbf{u}\mapsto\mathbf{v}$ alone, the model jointly
predicts $\boldsymbol{\kappa}$ and auxiliary quantities that co-determine the
coupled fixed point~\eqref{eq:coupled}.  Supervision is generated along the
OpenDSS daily march path, so labels reflect the trajectory-selected branch of the
multivalued map.  This aligns with Sandia's emphasis on state-aware or
time-correlated QSTS acceleration~\cite{deboever2017qsts,reno2017vq}.

\subsection{Optional physics-informed nodal loss}

When \texttt{--loss\_power\_balance\_weight} $> 0$, the trainer adds a mean-squared
residual on MV nodes:
\begin{align}
  r_P &= P_{\mathrm{inj}} - \Re\!\big\{ \mathbf{V}\,(\conj{\mathbf{Y}}\,\mathbf{V})^*\big\}, \\
  r_Q &= Q_{\mathrm{inj}} - \Im\!\big\{ \mathbf{V}\,(\conj{\mathbf{Y}}\,\mathbf{V})^*\big\},
\end{align}
with $\mathbf{Y}_{\mathrm{bus}}$ assembled from line $R/X$, \emph{predicted} LTC
tap ratios on regulator branches, and shunt $jB$ from \emph{predicted} capacitor
states; $P_{\mathrm{inj}},Q_{\mathrm{inj}}$ combine denormalized loads, meta-aux
PV totals, and predicted cap $Q$~\cite{gnn2_ieee8500_avc}.

\paragraph{Validation scope.}
\texttt{\_tmp\_pf\_physics\_validate.py} checks algebraic consistency on a
hand-built 4-bus network (residual $\approx 0$ at truth, grows under deliberate
perturbations, non-zero autograd).  This validates implementation fidelity of the
loss assembly, \textbf{not} end-to-end agreement with OpenDSS on the full 8500-node
feeder.  All preflight checks passed in the current repository snapshot.

% =============================================================================
\section{Evaluation Methodology}
\label{sec:eval}
% =============================================================================

\subsection{OpenDSS truth and metrics}

\texttt{nonunique.ipynb} cell~2 (\texttt{mode=da\_gps\_daily\_compare}) compares:
\begin{enumerate}
  \item \textbf{OpenDSS truth:} one compile, native \texttt{mode=daily}, sequential
        \texttt{Solve()} with warm-start carry-forward on profiles
        \texttt{load\_day\_004.csv} / \texttt{irr\_day\_004.csv}.
  \item \textbf{DA-GPS inference:} GNN-only forward passes on the same display
        grid (\texttt{step\_min}, \texttt{npts}), overlaying voltages, taps, and
        capacitor states on four monitor buses.
\end{enumerate}
Metrics include per-monitor and overall voltage MAE/RMSE (p.u.), regulator and
capacitor trajectories, control/PF iteration counts, and wall-clock timing.

\subsection{Representative results (notebook cell~3)}

On GPU, with \texttt{step\_min=5}~min ($n_{\mathrm{pts}}=288$):
\begin{itemize}
  \item OpenDSS mean \texttt{Solve()} time: $\approx 15.8$~ms/step (loop wall
        $\approx 4.56$~s).
  \item DA-GPS deployment mean: $\approx 4.9$~ms/display step (total deployment
        wall $\approx 1.40$~s).
  \item Wall speedup (OpenDSS loop / GNN deployment): $\approx 3.35\times$.
  \item Overall monitor voltage MAE: $\approx 4.3\times 10^{-3}$~p.u.; RMSE
        $\approx 5.2\times 10^{-3}$~p.u.
\end{itemize}
Coarser display grids ($10$--$60$~min) show similar or slightly higher speedups
($\approx 4\times$) with proportionally fewer forwards.

\subsection{Comparison to device-agnostic baselines}

Flat MLP sweeps on the same 8500-node dataset (\texttt{mlp\_sweep\_8500}) achieve
best test voltage MAE $\approx 0.25$~p.u.\ without graph structure or explicit
device-state heads---orders of magnitude worse than DA-GPS daily-compare MAE and
consistent with the failure mode of injection-only mappings under autonomous
controls~\cite{deboever2017qsts}.  Graph-based DA-GPS exploits topology and
multitask device supervision; timing comparisons in cell~3 are against OpenDSS
sequential QSTS, not against MLP inference latency.

% =============================================================================
\section{Discussion}
\label{sec:discussion}
% =============================================================================

\paragraph{Claims validated against the repository.}
\begin{itemize}
  \item \textbf{Sandia non-uniqueness / ML one-to-one mapping:} Supported by
        \cite{deboever2017qsts} (primary) and demonstrated empirically in cell~0.
  \item \textbf{IEEE~8500 unbalanced feeder vs OpenDSS baselines:} Supported;
        daily compare uses \texttt{Master-PV2MW-inv.dss} native daily QSTS.
  \item \textbf{PF + autonomous controls as closed-loop system:} Supported by
        \eqref{eq:coupled} and OpenDSS outer control loop
        documentation~\cite{gnn2_ieee8500_avc,opendss}.
  \item \textbf{Physics-informed nodal loss accuracy:} Implementation algebra
        validated synthetically; optional training term, not a substitute for
        OpenDSS verification on the full feeder.
\end{itemize}

\paragraph{Limitations.}
The surrogate is trained on OpenDSS trajectories and evaluated on paths that
include warm-start carry-forward; cold-start snapshot parity remains difficult
(cell~0 \texttt{snapshot\_cold}).  Physics loss uses predicted (not ground-truth)
controls inside $\mathbf{Y}_{\mathrm{bus}}$.  Speedups exclude one-time compile
and checkpoint load; wrapper overhead is reported separately in timing logs.

% =============================================================================
\section{Conclusion}
\label{sec:conclusion}
% =============================================================================

Autonomous voltage controls induce hysteresis and multiple valid coupled solutions
for identical injections, breaking naive QSTS surrogates that map
$\mathbf{u}\mapsto\mathbf{v}$~\cite{deboever2017qsts}.  Treating PF plus
controllers as a closed-loop fixed-point problem clarifies why state-aware,
multitask, graph-based training on sequential OpenDSS labels is appropriate for
the IEEE~8500 feeder.  DA-GPS achieves sub-$10^{-2}$~p.u.\ monitor MAE against
native daily QSTS while reducing deployment wall time by roughly $3$--$4\times$
on representative notebook configurations; optional nodal power-balance loss adds
physics consistency at the implementation level validated by synthetic tests.

% =============================================================================
\bibliographystyle{IEEEtran}
\bibliography{references}
% =============================================================================

\end{document}
